\subsection{Residual capacity and a common upper certificate}\label{sec:residual}
Fixing one bidder's complete mechanism leaves a constrained screening problem
for the other. This gives a useful interpretation of a capacity price.
For a jointly measurable capacity function $r:[0,1]^4\to[0,1]^2$, let
\[
\mathcal V(r)=\sup\{\operatorname{Rev}_2(M_2):
 M_2\text{ is pointwise DSIC/IR and }x_2\le r\}.
\]
The mechanisms in this definition remain jointly measurable, with integrable
payments. Grouping feasible pairs by their first component gives exactly
\[
\OPT=\sup_{M_1}\{\operatorname{Rev}_1(M_1)+\mathcal V(1-x_1)\}.
\]
This identity does not assert existence of a best response or invoke a
measurable selection of pointwise optimizers.

Write $a_i(t;w)=x_i(t,w)$ for bidder $i$'s allocation vector, with own
report $t$ and opponent report $w$. For a measurable, integrable payment
slice, put $R_i(w)=\int_{[0,1]^2}p_i(t,w)\,dt$, so that
$\operatorname{Rev}_i(M_i)=\int R_i(w)\,dw$. Such slices exist for almost
every $w$ in the stated jointly measurable class. On a fixed fiber define
\[
\begin{aligned}
 R(u,a)&=\int_{[0,1]^2}[t\cdot a(t)-u(t)]\,dt,\\
 V_w(s)&=\sup\{R(u,a):(u,a)\text{ admissible on the fiber},\ a\le s\}.
\end{aligned}
\]
Here admissible means $u\ge0$ is convex, $a(t)\in\partial u(t)$ and
$0\le a\le1$ at every report, with measurable allocation and integrable
payment $t\cdot a(t)-u(t)$. Thus $V_w$ is a conditional screening value;
$\mathcal V$ remains the value over complete jointly measurable mechanisms.
No equality between $\mathcal V$ and an integral of fiberwise suprema is assumed.

\begin{proposition}[Residual value and supporting measures]\label{prop:residual-value}
The functional $\mathcal V$ is nondecreasing and concave. On a fixed opponent
fiber $w$, suppose a finite nonnegative capacity-price measure $\pi$ satisfies
$R(u,a)\le\langle\pi,a\rangle$ for every admissible conditional mechanism,
and a feasible candidate at $r^*$ attains
$R^*=\langle\pi,r^*\rangle$. Then that candidate solves the full randomized
conditional problem and, for every residual for which the pairing is defined,
\[
V_w(r)\le V_w(r^*)+\langle\pi,r-r^*\rangle.
\]
\end{proposition}
\begin{proof}
Monotonicity is immediate. Mix any two admissible mechanisms using one
exogenous coin with constant probability $\theta$. All pointwise incentive,
IR and capacity inequalities are preserved, as are measurability and
integrability. Apply this to arbitrarily near-optimal competitors and pass
to the supremum. For the second assertion,
$R\le\langle\pi,a\rangle\le\langle\pi,r\rangle$, and equality at $r^*$
gives the claim. No own-report-dependent mixture is used.
\end{proof}


The common-certificate requirement is stronger than conditional optimality
at a single residual. On $\Omega=[0,1]^4$, suppose nonnegative finite
capacity measures $\Pi_j$ have well-defined pairings and satisfy
\[
 \operatorname{Rev}_i(M_i)\le\sum_j\int_\Omega x_{ij}\,d\Pi_j
\]
for every complete DSIC/IR mechanism of either bidder. Joint feasibility
then gives $\operatorname{Rev}_1(M_1)+\operatorname{Rev}_2(M_2)\le\sum_j\Pi_j(\Omega)$. Equality requires both
revenue inequalities and all priced capacity constraints to bind at the
same joint allocation. The construction below uses integrable densities,
so these pairings cover the measurable class of Section~\ref{sec:model}.
